% !TeX program = xelatex
\documentclass[11pt,a4paper]{article}
\usepackage{ctex}
\usepackage{geometry}
\geometry{left=2.5cm,right=2.5cm,top=2.5cm,bottom=2.5cm}
\usepackage{amsmath,amssymb}
\usepackage{booktabs}
\usepackage{array}
\usepackage{float}
\usepackage{setspace}
\setstretch{1.02}
\setlength{\parindent}{2em}
\setlength{\parskip}{0.2em}
\usepackage[colorlinks=true,linkcolor=black,citecolor=black]{hyperref}

\title{\bfseries AI工具使用详情说明}
\author{}
\date{}

\begin{document}
\maketitle

本参赛队在竞赛过程中使用人工智能工具作为辅助手段，核心建模、算法设计、求解与结果分析均由参赛队员独立完成。
本文档为论文参考文献声明所对应的详细说明，记录所用工具、使用目的、关键交互以及采纳与人工修改情况。
文内各条均可独立阅读，涉及的统计量与模型术语均为论文正文采用的标准表述。

\section*{一、所用AI工具清单}

所用AI工具及版本、开发机构、使用日期与主要用途汇总于表~\ref{tab:tools}。
表内“版本/型号”与“开发机构”信息用于论文参考文献的标注。

\begin{table}[H]
  \centering\footnotesize
  \setlength{\tabcolsep}{4pt}
  \renewcommand{\arraystretch}{1.3}
  \caption{所用AI工具清单}
  \label{tab:tools}
  \begin{tabular}{p{1.8cm}p{2.2cm}p{3.0cm}p{2.6cm}p{4.9cm}}
    \toprule
    工具名称 & 版本/型号 & 开发机构 & 使用日期 & 主要用途 \\
    \midrule
    Claude & Opus 5.5 & Anthropic & 2026-10-23 至 10-29 & 建模方法改进建议、LaTeX排版建议 \\
    DeepSeek & V4 Pro & 深度求索（DeepSeek） & 2026-10-23 至 10-29 & 建模方法讨论、部分求解器代码编写 \\
    DeepSeek & V4 Flash & 深度求索（DeepSeek） & 2026-10-23 至 10-29 & 文献与方法检索、绘图脚本片段 \\
    ChatGPT & Astra 6 & OpenAI & 2026-10-23 至 10-29 & 可视化方案与图表选取建议 \\
    ChatGPT & Sol 6 & OpenAI & 2026-10-23 至 10-29 & 数据预处理代码编写 \\
    GLM & 5.3 & 智谱AI（Zhipu AI） & 2026-10-23 至 10-29 & 文献检索、术语与表述核对 \\
    \bottomrule
  \end{tabular}
\end{table}

\section*{二、具体使用目的与应用环节}

AI工具仅用于建模方法改进建议、部分代码编写、可视化与排版建议、文献与方法检索四类环节。
工具不参与最终模型的选择、求解结果的判定与结论的作出，这些由参赛队员基于数据证据完成。

\textbf{（一）建模方法改进建议。}
用于问题一至问题三的方法设计讨论。
问题一的0到6级评分还原、问题二的关键化合物选择、问题三的化合物谱生成与跨风格迁移等环节，队员先提出方法思路，再请工具评估思路是否合理、指出潜在问题或补充可对照的方案。
工具给出的建议均须经过数据证据和交叉验证结果的检验，未经检验或被数据否定的建议不予采用。

\textbf{（二）部分代码编写。}
用于求解器、交叉验证、数据处理与绘图脚本的代码片段。
工具生成的代码由队员在本地环境运行并逐段核对，随机种子、折划分和数据口径按队内统一约定改写。

\textbf{（三）可视化与LaTeX排版建议。}
用于图表类型选取和论文排版。
工具建议图表类型（如热力图、小多图、森林图）与版式，最终图表由队内脚本从求解器产物生成，图表数值不取自工具输出。

\textbf{（四）文献与方法检索。}
用于检索同源文献、公开数据集与方法资料，核对术语与统计口径。
检索结果由队员核实原始出处后方可引用。

\section*{三、关键交互记录}

以下记录具有代表性的交互，提示词为大意，回复为要点。
全部回复均经队员核实、修改或舍弃。

\textbf{（1）问题一的等级还原与数据划分。}工具：Claude（Opus 5.5），建模辅助。
提示词大意：附件4只给出标准化后的连续评分，原始0到6分缺失；训练时如何避免同一品评团的信息同时进入训练集和测试集。
回复要点：利用并列最低值反推评分零点，按品评团分组做交叉验证，避免同团样本跨折。
采纳与修改：采纳“分组隔离”的思路，但将分组单位改为按25个日期组隔离并按风格分层；等级映射最终采用内层折外残差构造的预测分布中位数与分位数切点两种方案，经交叉验证选定，未采用回复中直接线性映射的建议。

\textbf{（2）问题二的冗余控制。}工具：DeepSeek（V4 Pro），建模辅助。
提示词大意：化合物之间强相关，直接按边际相关取前三会重复选到同一信息，如何避免冗余。
回复要点：按相关性构建共表达图并划分模块，同一模块至多选择一个，后续选择用偏相关剔除已选信息。
采纳与修改：采纳“图约束加偏相关前向选择”的框架；模块划分的具体方法、阈值与候选集合由队员对照多个阈值和模块纯度自行确定，并用模拟植入实验验证找回率后固定。

\textbf{（3）求解器与绘图脚本。}工具：DeepSeek（V4 Flash）、ChatGPT（Sol 6），代码编写。
提示词大意：请编写梯度提升树嵌套交叉验证的骨架，以及50维乘以3名次稳定性热力图的小多图布局。
回复要点：给出可运行的代码骨架与绘图布局。
采纳与修改：队员在本地运行并修正随机种子、折划分与标签，对输出数值逐项核对；所有图表数值来自队内求解器产物。

\textbf{（4）可视化方案选取。}工具：ChatGPT（Astra 6），可视化建议。
提示词大意：50个感官维度、每维3个名次的化合物入选频率如何在一张图中清晰呈现。
回复要点：建议用热力图加离散色阶，行按维度、列按名次，不显著项加斜线。
采纳与修改：采纳热力图方案，但色阶分档阈值、分组顺序与标注由队员按实际稳定性数值确定，图表由队内脚本生成。

\textbf{（5）文献与方法检索。}工具：GLM（5.3）、DeepSeek（V4 Flash），文献检索。
提示词大意：检索商用啤酒风味代谢物感官预测的同源文献及公开原始品评数据。
回复要点：定位到Nature Communications的啤酒风味论文及其Zenodo公开数据集。
采纳与修改：队员核对公开数据与赛题附件的一致性，数据来源与评分对齐由队员逐项验证，未直接采用工具给出的数据内容。

\section*{四、采纳内容与人工修改情况}

总体而言，AI工具的建议按“提出、核实、采纳或修改或舍弃”处理。
被采纳的是方法框架层面的思路（如分组隔离交叉验证、图约束前向选择、潜在高斯生成模型、热力图呈现）与代码骨架。
被修改或舍弃的是具体数值、参数、阈值与未经验证的结论。

所有公式由队员重新推导，所有数值由队内求解器在本地生成并经复核，所有图表由队内脚本绘制。
AI工具未直接生成论文中的任何最终数值、表格或图。
最终模型与结论由队员基于数据证据确定，队员对全部内容的正确性、可靠性与学术规范性负责。

\end{document}
